\section{Well-Typedness}

This section discusses the concept of well-typedness in the context of lambda calculus.
We first discuss types, typing judgements and typing rules. We then explore how to use these rules to show that a \lterm is well-typed.

Types of \lc are defined in \cref{def:types}, and are either $X$, which is the base type, or a function type, $A \fa A'$ with $A$ and $A'$ being types in \lc.
This formulation follows from Church's \emph{simple typed lambda calculus}, introduced in \cite{ALONZOCHURCH1940}.
In this thesis, we use the Curry-style notation where we write terms without explicit types.
This notation is formulated by Haskell B. Curry, and the notation used here is discussed in \cite{barendregt1992lambda}.

\begin{definition}
  \label{def:types}
  \textbf{Types of \boldmath${\lambda}$-calculus}
	\begin{align*}
    \begin{grammar}{
      \pr{$A$}{$X \gors A \fa A'$}
    }
    \end{grammar}
  \end{align*}
\end{definition}

A typing judgement is of the form $\judgement{\Gamma}{M : A}$, and it means that we can derive that $M$ has type $A$ from the assumptions in context $\Gamma$.
A context $\Gamma$ is a list of variable-type pairs, where each variable is associated with a type. 
A context is written as $\Gamma = x_1 : A_1, \ldots, x_n : A_n$, where each $x_i$ is a variable and each $A_i$ is a type.

The typing rules are defined in \cref{def:typing-rules}, and are from \cite{barendregt1992lambda}. The Id rule states that if $x$ has type $A$ in the context, we can conclude that $x$ has type $A$.
The \impliesi rule states that if we have a term $M$ of type $A'$ in the context $\Gamma, x : A$, we can conclude that from context $\Gamma$ it follows that $\lamb{x}{M}$ has type $A \fa A'$. 
The \impliese rule states that if we have a term $M$ of type $A \fa A'$ and a term $N$ of type $A$, we can conclude that the application $M N$ has type $A'$.

\begin{definition}
  \label{def:typing-rules}
  \textbf{Typing rules of \boldmath${\lambda}$-calculus}
  \begin{prooftree}
    \def\extraVskip{5pt}
    \AxiomC{\phantom{A}}
    \RightLabel{Id}
    \UnaryInfC{$\judgement{\Gamma, x : A}{x : A}$} % Conclusion
  \end{prooftree}
  
  \begin{prooftree}
    \def\extraVskip{5pt}
    \AxiomC{$\judgement{\Gamma, x : A}{M : A'}$}
    \RightLabel{\impliesi}
    \UnaryInfC{$\judgement{\Gamma}{\lamb{x}{M} : A \fa A'}$} % Conclusion
  \end{prooftree}
  
  \begin{prooftree}
    \def\extraVskip{5pt}
    \AxiomC{$\judgement{\Gamma}{M : A \fa A'}$}
    \AxiomC{$\judgement{\Gamma}{N : A}$}
    \RightLabel{\impliese}
    \BinaryInfC{$\judgement{\Gamma}{M N : A'}$} % Conclusion
  \end{prooftree}
\end{definition}

A \lterm is considered to be \emph{well-typed} if we can derive a typing judgement with the typing rules. 
The \lterms that are well-typed can be assigned a type. 
In \cref{ex:well-typedness}, we show that \lterm $(\bi(\lamb{x}{\bi x}))$ is well-typed by giving the full typing derivation.

\begin{example}
  \label{ex:well-typedness}
  \textbf{Well-typedness of \boldmath$(\bi(\lamb{x}{\bi x}))(\bi \bi)$}
  \small
  \begin{prooftree}  
    \def\extraVskip{5pt}
        \AxiomC{\phantom{A}}
        \RightLabel{Id}
        \UnaryInfC{$\judgement{y : A \fa A}{y : (A \fa A)}$}
      \RightLabel{\impliesi}
      \UnaryInfC{$\judgement{}{\lamb{y}{y} :  (A \fa A) \fa (A \fa A)}$}
  
            \AxiomC{\phantom{A}}
          \RightLabel{Id}
          \UnaryInfC{$\judgement{x : A , z : A}{z : A}$}
        \RightLabel{\impliesi}
        \UnaryInfC{$\judgement{x : A}{\lamb{z}{z} : A \fa A}$}
          \AxiomC{\phantom{A}}
        \RightLabel{Id}
        \UnaryInfC{$\judgement{x : A}{x : A}$}
        \RightLabel{\impliese}
        \BinaryInfC{$\judgement{x : A}{\bi x : A \fa A}$}
      \RightLabel{\impliesi}
      \UnaryInfC{$\judgement{}{\lamb{x}{\bi x} : A \fa A}$}
      \RightLabel{\impliese}
    \BinaryInfC{$\judgement{}{\bi (\lamb{x}{\bi x}) : A \fa A}$} % Conclusion
  \end{prooftree}
\end{example}

In contrast to the \lterm $(\bi(\lamb{x}{\bi x}))$, the lambda term $M M$ cannot be assigned a type. 
Since $M M$ is an application, we can only apply the \impliese rule, which would give a contradiction.
The \lterm $M$ cannot be both a function and the argument to the same function.
